\documentclass[a4paper,12pt]{article}
\usepackage[utf8]{inputenc}
\usepackage[ngerman]{babel}
\usepackage{amsmath,amssymb}
\usepackage{geometry}
\geometry{margin=2.5cm}

\title{Aufgabe 9 -- Eigenschaften von $Ax = b$}
\author{}
\date{}

\begin{document}
\maketitle

\section*{Fragestellung}

Sei $Ax = b$ ein lineares Gleichungssystem mit $A \in \mathbb{R}^{m \times n}$.
Sind die folgenden Situationen möglich? Wenn ja: einfaches Beispiel. Wenn nein: kurze Begründung.

\bigskip

\subsection*{(i) $m > n$ und das System hat eine eindeutige Lösung}

\textbf{Ja, das ist möglich.}

\textbf{Beispiel:} $m = 3$, $n = 2$:
\[
\begin{pmatrix} 1 & 0 \\ 0 & 1 \\ 1 & 1 \end{pmatrix}
\begin{pmatrix} x \\ y \end{pmatrix}
=
\begin{pmatrix} 1 \\ 2 \\ 3 \end{pmatrix}
\]

Das entspricht:
\[
\begin{cases}
x = 1 \\
y = 2 \\
x + y = 3
\end{cases}
\]

Die ersten zwei Gleichungen liefern eindeutig $x = 1, y = 2$.
Die dritte Gleichung ist damit auch erfüllt ($1 + 2 = 3$).

\textbf{Erklärung:} Bei $m > n$ gibt es mehr Gleichungen als Unbekannte. Wenn der Rang von $(A \mid b)$ gleich dem Rang von $A$ ist (konsistent) und der Rang gleich $n$ ist (alle Variablen sind Pivot-Variablen), dann gibt es eine eindeutige Lösung. Die überschüssigen Gleichungen sind dann redundant.

\bigskip

\subsection*{(ii) $m > n$ und das System hat unendlich viele Lösungen}

\textbf{Ja, das ist möglich.}

\textbf{Beispiel:} $m = 3$, $n = 2$:
\[
\begin{pmatrix} 1 & 1 \\ 2 & 2 \\ 3 & 3 \end{pmatrix}
\begin{pmatrix} x \\ y \end{pmatrix}
=
\begin{pmatrix} 1 \\ 2 \\ 3 \end{pmatrix}
\]

Das entspricht:
\[
\begin{cases}
x + y = 1 \\
2x + 2y = 2 \\
3x + 3y = 3
\end{cases}
\]

Alle drei Gleichungen sagen dasselbe aus: $x + y = 1$.
Der Rang ist 1, aber wir haben $n = 2$ Unbekannte.
Also gibt es eine freie Variable und unendlich viele Lösungen: $(x, y) = (1 - s,\; s)$ für $s \in \mathbb{R}$.

\textbf{Erklärung:} Obwohl $m > n$, kann der Rang von $A$ kleiner als $n$ sein, wenn die Gleichungen linear abhängig sind. Dann gibt es freie Variablen und somit unendlich viele Lösungen.

\bigskip

\subsection*{(iii) $m < n$ und das System hat eine eindeutige Lösung}

\textbf{Nein, das ist nicht möglich.}

\textbf{Begründung:}

Der Rang von $A$ kann höchstens $\min(m, n) = m$ sein (da $m < n$).
Also gilt:
\[
\text{Rang}(A) \leq m < n
\]

Falls das System konsistent ist, gibt es $n - \text{Rang}(A) \geq n - m > 0$ freie Variablen.
Das bedeutet: Wenn es überhaupt eine Lösung gibt, dann gibt es \textbf{unendlich viele}.

Falls das System inkonsistent ist, gibt es \textbf{keine} Lösung.

In beiden Fällen ist eine \textbf{eindeutige} Lösung unmöglich.

\end{document}
